% Modernised reading — fonds Alexandre Grothendieck, shelfmark 103, pages 1-17
% (the whole folder; page 7 is blank).
%
% This is an INTERPRETATION, not a transcription. It re-reads the pages in
% today's notation and vocabulary, states the hypotheses the manuscript leaves
% implicit, and drops the critical apparatus so that the mathematics reads
% straight through. Everything the brackets and underlines carried — a doubtful
% reading, a notation he used differently, a missing passage — moves into a
% footnote. It is held to being mathematically correct as it stands; where the
% pages are loose or mistaken, the true statement is what appears here and the
% footnote says what the page has. In any dispute of substance the
% transcription governs: whoever cites Grothendieck must cite batch-01.fr.tex.
%
% Scope: the folder holds one batch, pages 1-17, all transcribed, read whole
% before any of this was written; ONE file for the shelfmark.
%
% What the folder is. Three documents: (1) pp. 1-6, the typescript « Esquisse
% d'une théorie des Gr-Catégories » by Hoàng Xuân Sính, annotated in his hand
% (another copy of the typescript that folder 135 holds); (2) pp. 8-15, his
% letter to Deligne of 27.8.74, whose first four pages are the letter of
% folder 134-1, here four pages longer; (3) pp. 16-17, his typed letter of
% 23.6.1974 to Deligne, Giraud and Verdier about the jury of Sính's thesis.
% The spine: Sính's structure theorems on a point (Gr-categories classified by
% pi_0, pi_1, k in H^3; Picard categories by pi_0, pi_1, s), then their
% transposition to stacks over a topos in the letter, where each becomes the
% hypercohomology of a complex — the exact sequence (*) for Picard stacks,
% tau_{<=2}(RGamma(B_G mod X, N)[1]) for Gr-stacks — and where relative
% cohomology is needed and asked for. On the point, both letter statements
% collapse back to the typescript's theorems; the reading says so (ours).
%
% Order. Pages 12-13 are read after 14-15: page 12 refers to « B_e -> B_G
% plus haut », which only pages 14-15 introduce, and 13 -> 14 does not
% continue (13 ends « J'ai », 14 opens « Je te signale »). The transcription
% flags this; the reading follows the argument, and \pagerange follows it.
%
% Notation fixed for the whole document. C a Gr-category (today: 2-group),
% pi_0(C), pi_1(C) = Aut(1_C), k(C) in H^3(pi_0, pi_1); for a Picard category
% s(C) : pi_0 -> {}_2 pi_1. On a topos X: M, N abelian sheaves, E(M,N) =
% tau_{<=2} RHom(M,N), E'(M,N), P(M,N) = H^2(E'), sigma : P -> Hom(M, 2N).
% The structure sheaf is \mathcal O_X (his underlined O). B_G the classifying
% topos of a Group G of X, p_G : B_G -> X, q_G : X = B_e = (B_G)_{/P} -> B_G.
% RGamma(Y mod X, F) is the fibre of RGamma(Y, F) -> RGamma(X, q^*F).
%
% Substantive departures from the pages, each footnoted where it occurs:
% — p. 2: rho(X) is an automorphism of the GROUP pi_1(C) = Aut(1_C), not
%   « de 1_C ».
% — p. 1: in the H-space variant, E must be group-like (pi_0(E) a group) for
%   its objects to be invertible; supplied.
% — p. 2: « bi-Modules inversibles » read as bimodules locally free of rank
%   one on each side; Morita-invertible bimodules are more general.
% — p. 5: « isomorphisme de catégories » for the universal property of
%   S^{-1}C: only an equivalence is model-independent; stated so.
% — pp. 4-5: for S^{-1}C to be a Picard category in the sense of p. 1 (a
%   groupoid, by his own insertion) and for pi_1 = K^1, C must be the groupoid
%   of isomorphisms of the projective modules; supplied.
% — p. 9: « les invariants H^i [de E'] sont ceux de E en degré i+2 » is false
%   read literally; H^i(E') = H^i(E) for i <= 1 and (*) in degree 2 (as in
%   folder 134-1's reading).
% — p. 9: E'(M,N) is never defined; the definition over the sphere spectrum
%   (folder 134-1's reading) is given and marked as ours. The [1] on the
%   triangle's return arrow is ours.
% — p. 10: the obstruction lies in H^3(X, E(M,N)), a subgroup of
%   Ext^3(X; M, N).
% — p. 15: « champs de Picard ... épinglés par G, N » read as Gr-champs: G is
%   not assumed commutative.
% — p. 13: the cone problem is answered (ours): with functorial injective
%   resolutions the fibre is functorial at chain level and descends.
%
% Announced and not established: the construction of S^{-1}C (p. 5, « fort
% longue et fastidieuse »); K^0, K^1 of the envelope for projective modules
% (« on peut démontrer ») and for all additive categories (« sans doute »);
% the higher K^i through n-Picard categories (p. 5); a structure theory for
% the 2-categories of non-strict Picard and Gr-categories (p. 4, « dans les
% limbes »); the invariants of ⊕/⊗-categories (p. 6, a question); the triangle
% (T), heuristic, and the extension of Deligne's dictionary to strict
% 2-Picard stacks it assumes (pp. 8-9); the realisation of every sigma by a
% « universal » argument not given (p. 10); the descriptions of Gr-stacks by
% relative cohomology (pp. 14-15), asserted; a functorial RGamma(Y mod X, F)
% (pp. 12-13), asked; the question for Illusie (p. 13), whose answer begins
% on a leaf the folder does not hold. The report announced « ci-joint » on
% p. 16 is not in the folder. Page 17 withholds a private address, as the
% transcription does.
%
% Pass: Opus 5.5 (claude-opus-5-5), 2026-10-03 — first pass, unchecked against the pages by a human.
\documentclass[11pt,a4paper]{article}
\input{../preamble/grothendieck}

\folder{103}
\pages{1}{17}
\dating{1974}
\watermark{Édition de démonstration\\interprétation personnelle de l'œuvre}
\foldertitle{Courrier [Correspondance sur la thèse de Sinh] : copies d'article annoté (1974), lettres (1974) — lecture modernisée du dossier entier}

\begin{document}

\section*{Résumé}

\begin{resume}
Le dossier tient en trois pièces de l'année 1974, toutes trois autour de la
thèse de Hoàng Xuân Sính, mathématicienne vietnamienne qui y travaillait
depuis le séjour de Grothendieck au Vietnam en 1967 : un résumé de six
pages de cette thèse, tapé à la machine et corrigé de sa main ; une lettre à
Pierre Deligne écrite de Lodève à la fin d'août ; une lettre de juin à
Deligne, Giraud et Verdier, qui leur demande de siéger au jury.

Le sujet de la thèse se dit sans technique. Un groupe est un ensemble muni
d'une multiplication associative où tout élément a un inverse. Remplaçons
l'ensemble par une catégorie, dont les objets peuvent être isomorphes sans
être égaux, et l'égalité $(xy)z = x(yz)$ par un isomorphisme qu'il faut
choisir. On obtient ce que Sính appelle une Gr-catégorie, et qu'on appelle
aujourd'hui un $2$-groupe. Les lacets d'un espace en sont l'exemple de base :
mettre bout à bout deux lacets n'est associatif qu'à déformation près. La
question est de savoir ce qui, dans un tel objet, n'est pas déjà dit par les
deux groupes visibles --- les objets à isomorphisme près, et les
automorphismes de l'objet unité. La réponse est une seule classe de
cohomologie de degré $3$, qui mesure l'écart entre les isomorphismes
d'associativité choisis et l'identité, comme une classe de degré $2$ mesure
l'écart entre une extension de groupes et le produit. Si la multiplication est
de plus commutative à isomorphisme près, la réponse est encore plus petite :
un signe, ce que fait l'échange de deux copies d'un même objet. Le résumé
termine par l'application qui en fait le prix : la K-théorie algébrique d'un
anneau, ses deux premiers groupes $K^0$ et $K^1$, est exactement ce que voit
cette théorie lorsqu'on rend formellement inversibles les modules projectifs.

La lettre à Deligne fait passer tout cela d'un point à un espace --- plus
précisément à un topos, où les objets varient localement. Les mêmes
classifications reviennent, mais chacune s'enrichit d'un terme : ce qui, sur
un point, était un invariant isolé devient une suite exacte où apparaît un
groupe d'extensions qui, sur un point, est nul. Les théorèmes de Sính sont les
cas ponctuels de ce que la lettre écrit. Pour y parvenir dans le cas non
commutatif, Grothendieck a besoin d'une cohomologie relative à un morphisme de
topos, et il demande à Deligne si on sait la construire proprement.

Deux traits de manière se lisent sur ces pages. La lettre est faite
d'énoncés qu'il qualifie lui-même --- « reste heuristique », « sauf erreur »,
« ce qui pour moi ne fait guère de doute » --- et elle se termine par des
questions : est-ce connu, et de qui ? Et la dernière page du tapuscrit, corrigée de
sa main, annonce une fusion de l'algèbre homologique, de l'algèbre homotopique et de
l'algèbre catégorique en une « algèbre homologique non commutative », dont
ce qui porte alors ce nom « n'est que l'amorce » : c'est, neuf ans avant
\emph{Pursuing Stacks}, le programme de ce texte.

\keywords{2-group, categorical group, Sinh's theorem, Picard groupoid, Mac Lane obstruction, bitorsor, group completion, algebraic K-theory, rig category, homotopy hypothesis, Picard stack, relative cohomology, classifying topos, 2-gerbe}
\end{resume}

\section*{Le fil du dossier, et les conventions}

\begin{enumerate}
\item[1.] \emph{Le tapuscrit de Sính} (pages 1 à 6) : Gr-catégories et leur
classification par $(\pi_0, \pi_1, k)$ ; catégories de Picard et leur
classification par $(\pi_0, \pi_1, s)$ ; enveloppe de Picard d'une
$\otimes$-catégorie, $K^0$ et $K^1$ ; un programme. Page 7 est blanche.
\item[2.] \emph{La lettre du 27.8.74, première partie} (pages 8 à 11) : les
champs de Picard épinglés par deux faisceaux, le triangle $(T)$, la suite
exacte $(*)$, l'exemple de la K-théorie.
\item[3.] \emph{La même lettre, les Gr-champs} (pages 14 et 15) : leur
description par la cohomologie relative du topos classifiant $B_G$.
\item[4.] \emph{La même lettre, la cohomologie relative} (pages 12 et 13) :
deux constructions à la main, l'obstacle du cône, une question pour Illusie.
\item[5.] \emph{La lettre du 23.6.1974} (pages 16 et 17) : le jury de thèse.
\end{enumerate}

Les pages 12 et 13 sont lues ici après les pages 14 et 15. La page 12 renvoie
à « $B_e \to B_G$ plus haut », que seules les pages 14 et 15 introduisent ; la
page 13 s'achève sur « J'ai », que la page 14, ouverte par « Je te signale »,
ne continue pas. La suite 8--11, 14--15, 12--13 se lit sans rupture, et le
feuillet qui suivait la page 13 manque au dossier.\footnote{La transcription
garde l'ordre des archivistes et le signale ; on suit ici l'ordre de
l'argument.}

\emph{Le fil.} Le tapuscrit énonce sur un point deux théorèmes de structure :
une Gr-catégorie est déterminée par $\pi_0$, $\pi_1$ et une classe
$k \in H^3(\pi_0, \pi_1)$ ; une catégorie de Picard par $\pi_0$, $\pi_1$ et un
homomorphisme $s : \pi_0 \to {}_2\pi_1$. La lettre transpose l'un et l'autre
au-dessus d'un topos $X$, où les invariants deviennent des faisceaux et les
classifications des groupes d'hypercohomologie. Sur le topos ponctuel les deux
transpositions redonnent exactement les deux théorèmes, et c'est la meilleure
vérification qu'elles en soient bien les généralisations.\footnote{Cette
remarque est la nôtre ; elle est vérifiée aux sections correspondantes. La
lettre ne la fait pas, sans doute parce qu'elle va de soi entre les deux
correspondants.}

\emph{Conventions.} Une \emph{Gr-catégorie} est un $2$-groupe : un groupoïde
monoïdal dont tout objet est inversible. On garde le mot de Sính, qui est
encore en usage. On note $1_C$ l'objet unité,
$\pi_0(C)$ le groupe des classes d'objets, $\pi_1(C) = \mathrm{Aut}(1_C)$.
Les conditions AU, AC, ACU, AUC de Saavedra désignent une
$\otimes$-catégorie associative-unitaire, associative-commutative, etc. ; une
$\otimes$-catégorie AUC est aujourd'hui une catégorie monoïdale
symétrique.\footnote{N.~Saavedra Rivano, \emph{Catégories tannakiennes},
Lecture Notes in Math.\ 265, 1972 ; le tapuscrit renvoie à « Saavedra [ ] »
sans remplir les crochets. Toutes les références du tapuscrit sont ainsi des
crochets vides ; à deux endroits Grothendieck a écrit à la main « Deligne »
devant eux.} Pour un groupe abélien ou un faisceau abélien $N$, ${}_2N$ est le
sous-groupe (sous-faisceau) des éléments tués par $2$. Le faisceau structural
est noté $\mathcal{O}_X$.\footnote{Il le note par un $O$ souligné.} Pour un
morphisme de topos $q : X \to Y$ et un faisceau abélien $F$ sur $Y$,
$R\Gamma(Y \bmod X, F)$ est la fibre de $R\Gamma(Y, F) \to R\Gamma(X, q^*F)$,
de sorte que
\[ \cdots \to H^{i}(Y \bmod X, F) \to H^{i}(Y, F) \to H^{i}(X, q^{*}F) \to
H^{i+1}(Y \bmod X, F) \to \cdots \]

\pagerange{1}{2}
\section*{Les Gr-catégories et leurs exemples (pages 1 et 2)}

Le tapuscrit se présente comme le résumé de quelques résultats d'un travail
que l'auteur compte présenter comme thèse de doctorat.\footnote{Tapuscrit de
six pages, « Esquisse d'une théorie des Gr-Catégories, par Mme Hoang Xuan
Sinh ». C'est un autre exemplaire du tapuscrit que contient le dossier 135
(ses pages 39 à 44). Celui-ci est corrigé à la main : insertions en
interligne, mots tapés biffés, une note marginale, un paragraphe biffé ; un
contrôle de l'encre sur le fac-similé attribue ces additions à Grothendieck.
En haut à droite de la page 1, d'une autre main : « 9/74 Envoyé par
Grothendieck ». Les corrections qui changent le sens sont signalées en note
ci-dessous ; les autres, de pure rédaction, ne le sont pas.}

Soit $C$ une catégorie munie d'un produit $\otimes : C \times C \to C$,
associatif et unitaire à des isomorphismes donnés près, ces isomorphismes
satisfaisant les conditions de cohérence (pentagone et triangle). C'est une
\emph{Gr-catégorie} si $C$ est un groupoïde et si tout objet $X$ admet un
inverse $Y$, $X \otimes Y \simeq Y \otimes X \simeq 1_C$.\footnote{« c'est un
groupoïde, et si » est inséré à la main. La condition est indispensable :
sans elle, $\pi_1(C)$ ne mesurerait pas les automorphismes d'un objet
quelconque, et la classification qui suit serait fausse. Le nom moderne est
\emph{$2$-groupe}, ou \emph{groupe catégorique} ; « Gr-catégorie » est le
terme de Sính, resté en usage dans une partie de la littérature.}

\textbf{Exemples.}
\begin{enumerate}
\item[1.] Pour un espace pointé $(X, x)$, les lacets en $x$, avec pour
morphismes les classes d'homotopie d'homotopies de lacets et pour produit la
composition des lacets, associative à homotopie près. Variante : un espace de
Hopf $E$ associatif à homotopie près (en un sens assez fort) et
\emph{à inverse homotopique}, les objets étant ses points et les morphismes
les classes de chemins.\footnote{La condition que $\pi_0(E)$ soit un groupe
est ajoutée par nous ; sans elle les objets ne sont pas inversibles. Le
tapuscrit note que, lorsque $E$ admet un classifiant $X$, on retrouve
essentiellement l'exemple précédent. Une note marginale de sa main va dans
l'autre sens : le premier exemple est un cas particulier du second, en
prenant $E = \Omega(X, x)$, l'espace des lacets --- lecture douteuse quant à
la lettre, l'$\Omega$ portant peut-être un exposant. Les deux remarques sont
justes : à homotopie près, un espace de Hopf à inverse est l'espace des
lacets de son classifiant. Une ligne tapée après l'exemple est biffée au point
d'être illisible ; le mot « Exemple » est biffé devant 2 et 3, de sorte que
les trois forment une seule liste.}
\item[2.] Pour un faisceau abélien $F$ sur un espace ou un topos, les
$F$-torseurs avec la somme de Baer ; c'est même une catégorie de Picard
stricte. Les Modules inversibles sur un topos annelé en anneaux commutatifs
$(X, \mathcal{O}_X)$ en sont le cas $F = \mathcal{O}_X^{*}$.\footnote{Le
tapuscrit écrit « faisceau » ; la somme de Baer exige que $F$ soit
commutatif, ce que l'exemple suivant, où il dit « pas nécessairement
commutatif », confirme par contraste. « localement » est ajouté à la main
devant « annelé », et un mot tapé à sa suite est biffé.}
\item[3.] Pour une catégorie $A$, ou plus généralement un objet d'une
$2$-catégorie, les auto-équivalences de $A$ avec la composition. Si $F$ est un
faisceau en groupes, non nécessairement commutatif, et $A$ le champ des
$F$-torseurs à droite, ces auto-équivalences s'identifient aux
\emph{bitorseurs} sous $F$ : faisceaux sur lesquels $F$ opère à droite et à
gauche, de façon commutante, chacune des deux actions en faisant un torseur.
Pour $F = \mathcal{O}_X^{*}$, $\mathcal{O}_X$ faisceau d'anneaux non
nécessairement commutatif, ce sont les $\mathcal{O}_X$-bimodules localement
libres de rang un de chaque côté.\footnote{Le tapuscrit dit « bi-Modules
inversibles », et Grothendieck a mis « inversibles » entre guillemets à la
main. La précaution est fondée : un bimodule inversible au sens de Morita
n'est pas en général libre de rang un --- les bimodules qui réalisent
l'équivalence entre un anneau et un anneau de matrices sur lui ne le sont
pas. Ceux qui correspondent aux bitorseurs sous $\mathcal{O}_X^{*}$ sont
exactement ceux qui admettent localement une base $e$ à gauche ; alors
$e\,a = \sigma(a)\,e$ pour un automorphisme $\sigma$ de l'anneau, et $e$ est
aussi une base à droite.}
\end{enumerate}

\pagerange{2}{3}
\section*{Le théorème de structure : $\pi_0$, $\pi_1$, $k$ (pages 2 et 3)}

À une Gr-catégorie $C$ on associe :
\begin{enumerate}
\item[a)] le groupe $\pi_0(C)$ des classes d'isomorphisme d'objets ;
\item[b)] le groupe $\pi_1(C) = \mathrm{Aut}(1_C)$, qui est commutatif ;
\item[c)] une action de $\pi_0(C)$ sur $\pi_1(C)$. Pour tout objet $X$, les
deux applications $u \mapsto u \otimes \mathrm{id}_X$ et
$u \mapsto \mathrm{id}_X \otimes u$, composées avec les isomorphismes
d'unité, sont des isomorphismes de groupes
$\gamma_X, \delta_X : \mathrm{Aut}(1_C) \to \mathrm{Aut}(X)$ ;
l'automorphisme $\rho(X) = \delta_X^{-1} \circ \gamma_X$ du groupe
$\pi_1(C)$ ne dépend que de la classe de $X$, et fait de $\pi_1(C)$ un
$\pi_0(C)$-module.\footnote{Le tapuscrit appelle $\rho(X)$ « l'automorphisme
de $1_C$ déduit des deux isomorphismes » ; c'est un automorphisme du
\emph{groupe} $\mathrm{Aut}(1_C)$, non de l'objet $1_C$, et c'est ainsi que
la suite l'emploie. La commutativité de $\pi_1(C)$ est l'argument
d'Eckmann--Hilton.}
\end{enumerate}

Choisissons pour chaque $a \in \pi_0$ un représentant $L_a$, avec $L_1 = 1_C$,
et pour chaque couple un isomorphisme
$\varphi_{a,b} : L_a \otimes L_b \simeq L_{ab}$. L'isomorphisme d'associativité
$(L_a \otimes L_b) \otimes L_c \simeq L_a \otimes (L_b \otimes L_c)$, lu à
travers $\varphi_{a,b}$, $\varphi_{ab,c}$, $\varphi_{b,c}$, $\varphi_{a,bc}$,
devient un automorphisme de $L_{abc}$, c'est-à-dire un élément
$f(a,b,c) \in \pi_1$. Le pentagone dit que $f$ est un $3$-cocycle de $\pi_0$ à
valeurs dans le $\pi_0$-module $\pi_1$ ;\footnote{« (axiome du pentagone) »
est ajouté à la main.} changer les $\varphi$ par une $2$-cochaîne $g$ change
$f$ en $f + dg$. D'où une classe bien définie
\[ k(C) \in H^3\bigl(\pi_0(C), \pi_1(C)\bigr) . \]

\begin{quote}
\textbf{Théorème (Sính).} Une Gr-catégorie est déterminée, à
$\otimes$-équivalence près, par le groupe $\pi_0$, le $\pi_0$-module $\pi_1$
et la classe $k$ ; toute classe se réalise.
\end{quote}

Précisément : deux Gr-catégories sont équivalentes si et seulement s'il
existe des isomorphismes $\alpha : \pi_0 \simeq \pi'_0$ et
$\beta : \pi_1 \simeq \pi'_1$, compatibles aux actions, qui transportent $k$
sur $k'$.\footnote{La formulation précise est la nôtre ; le tapuscrit dit
« la classification se fait précisément en termes de » ces données. La
classe d'une Gr-catégorie dont $\pi_0$ et $\pi_1$ sont identifiés à des
groupes fixés est bien définie ; sans identification, elle ne l'est qu'à
l'action des automorphismes près. Le tapuscrit ajoute que la loi de groupe de
$H^3$ s'interprète « à la Baer » par des opérations sur les Gr-catégories.}
C'est aujourd'hui le théorème de Sính, et la version catégorique de la
classification des $2$-types d'homotopie connexes par leur premier invariant
de Postnikov.

\emph{Cas particuliers.} Dans l'exemple~1, on doit retrouver le premier
invariant de Postnikov de l'espace, dans $H^3(\pi_1(X), \pi_2(X))$, avec une
interprétation géométrique nouvelle.\footnote{« doit » est inséré à la main :
il écrit « on doit retrouver » là où le tapuscrit affirmait « on retrouve ».
La correction est de prudence, l'énoncé étant vrai.} Dans l'exemple~2,
$\pi_0 = H^1(X,F)$, $\pi_1 = H^0(X,F)$ et $k = 0$, puisqu'une catégorie de
Picard stricte est déterminée par $\pi_0$ et $\pi_1$ (section
suivante).\footnote{Le tapuscrit renvoie ici à « [ ] », et Grothendieck a
écrit « Deligne » devant le crochet.} Dans l'exemple~3, $\pi_0$ est le groupe
des classes d'auto-équivalences de $A$, $\pi_1$ le groupe des automorphismes
du foncteur identique, et $k$ semble nouveau ; pour les bimodules,
$\pi_0$ s'interprète par un « groupoïde de Brandt » et $\pi_1$
est le groupe des unités centrales de $\mathcal{O}_X$. Lorsque $A$ est la
catégorie des torseurs sous un groupe ordinaire $F$, on trouve
$\pi_0 = \mathrm{Out}(F)$, $\pi_1 = Z(F)$, et $k$ est la classe de Mac Lane :
pour un noyau abstrait $\psi : Q \to \mathrm{Out}(F)$, l'obstruction de Mac
Lane à réaliser $\psi$ par une extension est l'image réciproque
$\psi^{*}k \in H^3(Q, Z(F))$.\footnote{La formulation par image réciproque est
la nôtre ; le tapuscrit dit que $k$ « n'est autre que le classique invariant
de Mac-Lane ». Il termine en souhaitant un lien \emph{direct} et géométrique
entre la théorie des extensions de Mac Lane et celle des Gr-catégories et de
leurs « splittages ». « développer » est écrit au-dessus d'un mot tapé biffé
(« établir », sans certitude), « géométrique » au-dessus de « direct », qui
n'est pas biffé : on garde les deux.}

\pagerange{3}{4}
\section*{Les catégories de Picard (pages 3 et 4)}

Une \emph{catégorie de Picard} est une $\otimes$-catégorie ACU qui est une
Gr-catégorie pour sa structure AU : aujourd'hui, un \emph{groupoïde de
Picard}, ou $2$-groupe symétrique. Elle est \emph{stricte} si la symétrie
$c_{X,X}$ de $X \otimes X$ est l'identité pour tout $X$. Les catégories de
Picard strictes sont celles que Deligne étudie, dans le cadre plus large des
champs de Picard.\footnote{SGA 4, exposé XVIII, § 1.4, paru en 1973.}

Pour une catégorie de Picard, $\pi_0$ est commutatif et opère trivialement
sur $\pi_1$. Une équivalence de Gr-catégories entre deux catégories de Picard
n'est pas forcément compatible aux symétries, de sorte que la classification
au sens Picard est plus fine. La symétrie $c_{X,X}$, lue comme un élément de
$\pi_1$, définit un \emph{homomorphisme}
\[ s(C) : \pi_0 \longrightarrow {}_2(\pi_1) . \]
C'est un homomorphisme parce que
$c_{X \otimes Y, X \otimes Y} = c_{X,X}\, c_{Y,Y}\, c_{X,Y}\, c_{Y,X}$ et que
les deux derniers facteurs sont inverses l'un de l'autre ; il est à valeurs
dans ${}_2\pi_1$ parce que $c_{X,X}^{\,2} = \mathrm{id}$.\footnote{« homomorphisme »
remplace à la main « application », tapé et biffé. La justification est
ajoutée par nous ; le tapuscrit n'en donne pas. Les indices de la formule
sont retouchés à la main sur le tapuscrit, qui les avait tapés autrement ; on
la donne telle que la ligne suivante la fixe.}

\begin{quote}
\textbf{Théorème (Sính).} Une catégorie de Picard est déterminée, à
équivalence de catégories de Picard \emph{non canonique} près, par les groupes
commutatifs $\pi_0$, $\pi_1$ et l'homomorphisme $s : \pi_0 \to {}_2\pi_1$.
Elle est stricte si et seulement si $s = 0$ ; une catégorie de Picard stricte
est donc déterminée par $\pi_0$ et $\pi_1$ seuls.\footnote{« non canonique »
est ajouté à la main. Le dernier énoncé est attribué par le tapuscrit à
« [ ] » --- c'est-à-dire à Deligne, comme le précisent les insertions
voisines.}
\end{quote}

\emph{Remarques.} Le théorème de Deligne décrit la $2$-catégorie des catégories
de Picard strictes par les complexes de groupes abéliens $K$ tels que
$H_i(K) = 0$ pour $i \neq 0, 1$ : la catégorie obtenue en identifiant les
$1$-flèches isomorphes est équivalente à la sous-catégorie pleine de la
catégorie dérivée $D(\mathrm{Ab})$ formée de ces complexes, avec
$H_0 = \pi_0$ et $H_1 = \pi_1$.\footnote{Le tapuscrit dit « interpréter la
$2$-catégorie […] en termes de la sous-catégorie pleine de la catégorie
dérivée » ; la précision sur les $1$-flèches est la nôtre, une catégorie
dérivée ne voyant pas les $2$-flèches. Les $H$ de cette remarque portent
chacun une marque à la main sur l'indice, dont l'intention n'est pas établie
--- peut-être un passage en exposant, qui donnerait l'indexation
cohomologique $D^{[-1,0]}$ ; un « alors » manuscrit court sous la formule,
sans signe d'insertion lisible.} Le théorème précédent, « assez grossier »,
se réduit alors à ceci : un tel complexe est déterminé à isomorphisme non
unique près par son $H_0$ et son $H_1$, ce qui tient à ce que
$\mathrm{Ext}^2_{\mathbf{Z}}$ est nul. Il resterait à faire une théorie de
structure de la $2$-catégorie des catégories de Picard non strictes, et de
celle des Gr-catégories, et à transposer tout cela aux champs, comme Deligne
l'a fait pour les catégories de Picard strictes.\footnote{« développements »
remplace « résultats », biffé ; le mot qui précède « étudiés » est biffé et
illisible. Le tapuscrit dit ces développements « encore dans les limbes ». La
transposition aux champs est exactement ce que fait la lettre des pages 8 à
15.}

\pagerange{4}{6}
\section*{L'enveloppe de Picard, $K^0$ et $K^1$ (pages 4 à 6)}

Soient $C$ une $\otimes$-catégorie AUC et $S$ une partie de $\mathrm{Ob}\,C$.
On cherche un $\otimes$-foncteur AUC universel $C \to S^{-1}C$ rendant
inversibles les objets de $S$ : pour toute $\otimes$-catégorie AUC $G$, la
restriction
\[ \underline{\mathrm{Hom}}_{\otimes \mathrm{AUC}}(S^{-1}C, G) \longrightarrow
   \underline{\mathrm{Hom}}^{S}_{\otimes \mathrm{AUC}}(C, G) \]
vers les $\otimes$-foncteurs qui inversent les objets de $S$ doit être une
équivalence de catégories. Le tapuscrit affirme que le problème a toujours une
solution, la \emph{$\otimes$-catégorie de fractions}, et n'en donne pas la
construction, « assez naturelle » mais dont l'explicitation est « fort longue
et fastidieuse ».\footnote{Le tapuscrit dit « isomorphisme de catégories ».
Pour une construction particulière il peut l'être, mais la propriété
universelle n'en fixe la solution qu'à équivalence près, et c'est
l'équivalence qui est la forme invariante de l'énoncé. La construction est
donnée dans le manuscrit anglais du dossier 135. À la page 4, « catégorie »
est complété en « $\otimes$-catégorie AUC $G$ » par une insertion manuscrite.}

Lorsque $S = \mathrm{Ob}\,C$, et que $C$ est un groupoïde, $S^{-1}C$ est une
catégorie de Picard, en général non stricte : l'\emph{enveloppe de Picard} de
$C$, dont les invariants sont notés
\[ K^0(C) = \pi_0(S^{-1}C), \qquad K^1(C) = \pi_1(S^{-1}C) . \]
Si $C$ est le groupoïde des modules projectifs de type fini sur un anneau $A$
et de leurs isomorphismes, avec la somme directe, ce sont les groupes
$K^0(A)$ de Grothendieck et $K^1(A)$ de Whitehead et Bass. Le tapuscrit tient
pour probable qu'il en va de même pour toute catégorie additive munie de la
somme.\footnote{L'hypothèse de groupoïde est ajoutée par nous ; le tapuscrit
dit « la catégorie des modules projectifs de type fini ». Elle est requise
par sa propre définition : une catégorie de Picard est une Gr-catégorie, donc,
depuis l'insertion de la page 1, un groupoïde. On sait aujourd'hui que
l'énoncé est vrai pour le groupoïde des isomorphismes de toute catégorie
additive, avec $K^1$ au sens de Bass pour la somme directe : l'enveloppe de
Picard est le tronqué en degré $1$ de la complétion en groupe, celle du
$S^{-1}S$ de Quillen, et Deligne la retrouve en 1987 comme catégorie des
« objets virtuels » (\emph{Le déterminant de la cohomologie}). « avec
l'opération de somme » est ajouté à la main, et « on peut démontrer que »
reste tapé sur cet exemplaire.}

Cette enveloppe n'est pratiquement jamais stricte. L'invariant $s$ y est
explicite : la symétrie de $P \oplus P$ est l'échange des deux facteurs, dont
la classe dans $K^1(A)$ est celle de $-\mathrm{id}_P$, de sorte que
$s([P]) = [-\mathrm{id}_P]$.\footnote{Le calcul est le nôtre : la matrice
d'échange diffère de $\mathrm{diag}(1,-1)$ par une matrice élémentaire. Il
précise le « pratiquement » du tapuscrit : l'enveloppe est stricte lorsque
$2 = 0$ dans $A$, puisqu'alors $-1 = 1$, et ne l'est pas, par exemple, pour
$A = \mathbf{Z}$, où $K^1 = \{\pm 1\}$ et où $s$ est la réduction modulo $2$
du rang.} D'où une conséquence négative : il est peu probable qu'on construise
dans $D(\mathrm{Ab})$ un complexe \emph{canonique} $K$ avec
$H_0(K) = K^0$ et $H_1(K) = K^1$, car une telle construction donnerait, par
la théorie de Deligne, une catégorie de Picard stricte canonique ayant ces
invariants.\footnote{L'argument est exact et il nomme l'obstruction en
termes d'aujourd'hui : le spectre de K-théorie n'est pas un module sur
$H\mathbf{Z}$, son invariant $s$ étant la multiplication par l'élément
$\eta$ de la sphère, et un complexe ne fournit que des
$H\mathbf{Z}$-modules. Un complexe non canonique ayant ces homologies existe
toujours, $\mathrm{Ext}^2_{\mathbf{Z}}$ étant nul ; c'est la canonicité qui
est en cause.}

\subsection*{Le programme (pages 5 et 6)}

Les $K^i$ supérieurs devraient s'interpréter de même par des
$n$-catégories de Picard et la construction de $n$-catégories de Picard
enveloppantes, dont les invariants homotopiques $\pi_i$ seraient les
$K^i$.\footnote{C'est ce qui s'est fait, sous le nom de complétion en
groupe : la K-théorie d'une catégorie monoïdale symétrique est un spectre
connectif, c'est-à-dire un $\infty$-groupoïde de Picard, et ses groupes
d'homotopie sont les $K^i$ de Quillen (1972--1973 ; G.~Segal, \emph{Categories
and cohomology theories}, 1974). Le tapuscrit ne cite pas ces travaux, et
rien n'indique qu'il les ait connus. Au bas de la page 5, un paragraphe tapé
est biffé de grands traits : il proposait une autre direction, la
$\otimes$-catégorie de fractions d'une $\otimes$-catégorie AU, sans
commutativité, et s'interrompt au milieu d'un mot.} De tels développements
obligeraient à tirer au clair les relations entre $n$-catégories et
$\infty$-catégories d'une part, ensembles simpliciaux et espaces d'autre part ;
entre $n$-Gr-catégories et groupes simpliciaux ou topologiques ; entre
$n$-catégories de Picard strictes et complexes de chaînes tronqués en
dimension $n$. On doit s'attendre à une fusion de l'algèbre homologique, de
l'algèbre homotopique et de l'algèbre catégorique dans une « algèbre
homologique non commutative » : « ce qui porte actuellement ce nom n'est que
l'amorce de ce qui doit venir ».\footnote{« doit » remplace à la main « reste
à ». La première de ces relations est ce qu'on appelle aujourd'hui
l'hypothèse d'homotopie, que \emph{Pursuing Stacks} (1983) place au centre ;
la dernière est la correspondance de Dold--Kan (1958), à laquelle les
catégories de Picard strictes de Deligne donnent un premier étage
catégorique.}

Dernière question, plus terre à terre. Si $A$ est un anneau commutatif, ou
plus généralement une catégorie additive munie d'un produit $\otimes$ ACU
biadditif, l'enveloppe de Picard porte, outre la somme $\oplus$, un produit
$\otimes$ venu du produit tensoriel, lié à $\oplus$ par des isomorphismes de
distributivité
\[ X \otimes (Y \oplus Z) \simeq (X \otimes Y) \oplus (X \otimes Z) . \]
Il faudrait étudier systématiquement les catégories munies de deux telles
opérations, de Picard pour $\oplus$ ; $\pi_0$ y est un anneau commutatif
unitaire et $\pi_1$ un module sur lui. Quels autres invariants les
caractérisent à équivalence près ?\footnote{« ACU », « par », « ordinaire »
et le début de la phrase, « suivante. Lorsque », sont des insertions
manuscrites. Le tapuscrit appelle $A$ à la fois l'anneau et la catégorie
munie des deux opérations. Ces catégories sont aujourd'hui les catégories
bimonoïdales ou « rig-catégories » ; la cohérence de la distributivité est
due à M.~Laplaza (1972), que le tapuscrit ne cite pas.}

\pagerange{8}{11}
\section*{Lettre à Deligne du 27.8.74 : les champs de Picard épinglés (pages 8 à 11)}

La lettre est écrite de Lodève le 27 août 1974.\footnote{La date est pâle sur
cet exemplaire : « 7.8.74 » se lit nettement, le 2 qui précède est presque
effacé ; il est lisible sur l'autre exemplaire de la lettre, celui du dossier
134-1, et vérifié sur le fac-similé. Les pages 8 à 11 sont la lettre de
quatre pages que contient ce dossier-là ; l'exemplaire présent se poursuit
quatre pages de plus. Dans la marge de gauche de la page 8, à la verticale,
d'une main que rien ne distingue de celle de la lettre : « Je n'ai pas
retrouvé la lettre sur les Gr-champs que tu cites ». La page ne dit pas qui
l'a écrite ; on notera que la page 14 s'ouvre sur « Je te signale que j'ai
réfléchi aux Gr-champs ».} Empêché peut-être d'assurer un cours de premier
cycle, Grothendieck envisage un petit séminaire d'algèbre sur « les fourbis
de Mme Sinh », éventuellement transposés aux champs ; c'est en faisant la
transposition qu'il « tombe sur le truc suivant ».\footnote{« par une jambe »,
lu avec doute, donne la raison de l'empêchement.} La lecture qui suit est
celle du dossier 134-1, résumée ; on s'y reporte pour les
démonstrations.\footnote{Les notations et les définitions supplées sont les
mêmes dans les deux lectures, pour qu'on puisse passer de l'une à l'autre.}

Soient $X$ un topos et $M$, $N$ deux faisceaux abéliens sur $X$. Un champ de
Picard \emph{épinglé} par $M$, $N$ est un champ de Picard $P$ sur $X$ muni
d'isomorphismes $\pi_0(P) \simeq M$, $\pi_1(P) \simeq N$. Son invariant de
symétrie $\sigma(P) : M \to {}_2N$ est le $s$ du tapuscrit, faisceautisé :
la symétrie de $L \otimes L$ pour une section locale $L$, lue dans
${}_2N$.\footnote{« s'obtient en » et « $L$ » sont insérés en interligne,
« champs » remplace « catégories », biffé ; un symbole est biffé devant
$L \otimes L$. L'exemplaire du dossier 134-1 porte les mêmes corrections.}
Posons
\[ E(M,N) = \tau_{\leq 2}\, R\underline{\mathrm{Hom}}(M,N), \qquad
   \underline{H}^i\bigl(E(M,N)\bigr) = \underline{\mathrm{Ext}}^i(M,N)
   \quad (0 \leq i \leq 2), \]
nul en dehors de $[0,2]$ ; son $\underline{H}^2$ classe, par le théorème de
Deligne, les champs de Picard \emph{stricts} épinglés par $M$, $N$. La lettre
pose un second complexe $E'(M,N)$ qui ferait de même pour les champs de
Picard non nécessairement stricts, et qu'elle ne définit pas ; on le
définit comme le complexe tronqué des morphismes de $HM$ dans $HN$ calculés
au-dessus de la sphère $\mathbb{S}$ et non plus de $\mathbf{Z}$,
\[ E'(M,N) = \tau_{\leq 2}\, R\underline{\mathrm{Hom}}_{\mathbb{S}}(HM, HN) .
\]
Cette définition est la nôtre.\footnote{Elle est celle de la lecture du
dossier 134-1, et anachronique quant au langage. Ce qu'elle formalise est ce
que la lettre dit de $E'$ : il « exprime le $2$-champ de Picard strict formé
des $1$-champs de Picard (pas nécessairement stricts) épinglés par $M$,
$N$ », « en admettant que la théorie pour les $1$-champs de Picard stricts
s'étend aux $2$-champs de Picard stricts (ce qui pour moi ne fait guère de
doute) ». « épinglés par $M$, $N$ » et « stricts » sont insérés en
interligne.} Le morphisme d'oubli $E \to E'$ s'insère alors dans un triangle
distingué $(T)$ :

\begin{tikzcd}
& \underline{\mathrm{Hom}}(M, {}_2 N)[-2] \arrow[dl, "{[1]}"'] & \\
E(M,N) \arrow[rr] & & E'(M,N) \arrow[ul]
\end{tikzcd}

dont le troisième terme vient de ce que l'homologie entière de $H\mathbf{Z}$
vaut $\mathbf{Z}$, $0$, $\mathbf{Z}/2$ en degrés $0$, $1$, $2$.\footnote{La page
dessine le triangle sans marque de décalage ; le « $[1]$ » est le nôtre. Le
sens des flèches est celui de la page, et c'est lui qui donne $(*)$.} La
suite de cohomologie de $(T)$ donne
$\underline{H}^i(E') = \underline{\mathrm{Ext}}^i(M,N)$ pour $i \leq 1$ et,
en degré $2$, la suite exacte
\[ (*) \qquad 0 \longrightarrow \underline{\mathrm{Ext}}^2(M,N)
   \longrightarrow P(M,N) \xrightarrow{\ \sigma\ }
   \underline{\mathrm{Hom}}(M, {}_2 N) \longrightarrow 0 , \]
où $P(M,N) = \underline{H}^2(E'(M,N))$ est le faisceau des classes, à équivalence près, de champs de
Picard épinglés par $M$, $N$.\footnote{La page dit que les invariants
$\underline{H}^i$ de $E'$ « sont ceux de $E$ en degré $i+2$ » ; pris à la
lettre c'est faux, le membre de droite étant nul pour $i > 0$, et l'énoncé
exact est celui qu'on donne. La lettre ajoute que $(*)$ se construit « en
tous cas canoniquement à la main », le terme médian étant le faisceau des
« classes » --- mot lu avec doute --- à équivalence près de champs de Picard
épinglés ; la parenthèse qui ferme « à la main » n'a pas d'ouvrante sur la
page.}

Sur le topos ponctuel, $\mathrm{Ext}^2_{\mathbf{Z}}$ est nul et $(*)$ se
réduit à $P(M,N) \simeq \mathrm{Hom}(M, {}_2N)$ : c'est exactement le
théorème de Sính de la page 4. La suite $(*)$ en est donc la forme
faisceautique, et le terme $\underline{\mathrm{Ext}}^2$ est ce qu'ajoute le
passage d'un point à un topos.\footnote{Cette lecture est la nôtre.}

\emph{Tout signe se réalise.} Grothendieck dit savoir prouver, « sauf
erreur », que tout homomorphisme $M \to {}_2N$ provient d'un champ de Picard
épinglé par $M$, $N$ : l'obstruction est a priori de degré $3$, et un argument
« universel », qu'il ne donne pas, prouve qu'elle est nulle.\footnote{La page
place l'obstruction dans $\mathrm{Ext}^3(X; M, N)$. Le groupe où elle vit
est $\mathbb{H}^3(X, E(M,N))$, qui s'injecte dans $\mathrm{Ext}^3(X; M, N)$ ;
la précision est la nôtre, et la lecture du dossier 134-1 propose une
reconstruction de l'argument, par réduction au cas où $M$ est libre.} Ainsi
toute section de $\underline{\mathrm{Hom}}(M, {}_2N)$ sur un objet $U$ se
relève dans $P(M,N)$, et même dans $\mathbb{H}^2(U, E'(M,N))$ : l'extension
$(*)$ a une section ensembliste, elle est « bien proche d'être splittée »,
sans être dite scindée.

\emph{L'exemple.} Pour un anneau $A$ sur $X$, soient $M$, $N$ les faisceaux
$K^0$ et $K^1$ associés au champ additif des $A$-Modules projectifs de type
fini. L'enveloppe de Picard des pages 4 et 5, faite localement, est un champ de
Picard épinglé par $M$, $N$, et donne une section canonique de $P(M,N)$. Son
invariant $\sigma$ envoie la classe de $P$ sur celle de $-\mathrm{id}_P$,
comme au calcul de la page 5.\footnote{L'identification de $\sigma$ est la nôtre ; la lettre
s'arrête à l'existence de la section. Un « de » est biffé avant
« canonique ».} Tout cela a « les fonctorialités évidentes en $M$, $N$,
$X$ ».

\emph{La question.} Le triangle $(T)$ et la suite $(*)$ sont-ils connus « par
les compétences (Quillen, Breen, Illusie …) » ? Ont-ils des variantes
supérieures, qu'on obtiendrait peut-être par des $n$-champs de Picard non
nécessairement stricts ?\footnote{« compétences » est vérifié sur le
fac-similé, ici comme dans le dossier 134-1. En termes d'aujourd'hui, les
variantes supérieures sont les étages de la comparaison entre
$R\underline{\mathrm{Hom}}_{\mathbb{S}}$ et
$R\underline{\mathrm{Hom}}_{H\mathbf{Z}}$, gouvernée par l'homologie entière
de $H\mathbf{Z}$ ; voir la lecture du dossier 134-1. La lettre demande si la
chose est connue, ce qui est le contraire d'une revendication.}

\pagerange{14}{15}
\section*{Les Gr-champs épinglés (pages 14 et 15)}

Soient $G$ un Groupe sur $X$, non nécessairement commutatif, et $N$ un
$G$-Module. Les Gr-champs sur $X$ épinglés par $(G, N)$ --- c'est-à-dire
munis d'isomorphismes $\pi_0 \simeq G$, $\pi_1 \simeq N$ compatibles à
l'action --- forment une $2$-catégorie, et même une $2$-catégorie de Picard
stricte grâce à l'opération de Baer.\footnote{« épinglé par $(G,N)$ » est
écrit dans la marge de gauche avec un trait de renvoi, « (de cochaînes) »
en interligne.} Soient $B_G$ le topos classifiant de $G$ (les objets de $X$
munis d'une action de $G$), $p_G : B_G \to X$ sa projection, $P$ le torseur
universel ($G$ agissant sur lui-même) et
$q_G : X \simeq B_e \simeq (B_G)_{/P} \to B_G$. Le complexe de cochaînes
tronqué qui correspond à cette $2$-catégorie de Picard stricte est
\[ \tau_{\leq 2}\bigl( R\Gamma(B_G \bmod X, N)[1] \bigr) , \]
la cohomologie de $R\Gamma(B_G \bmod X, N)$ commençant en degré $1$.
Géométriquement, un Gr-champ épinglé par $(G, N)$ est essentiellement une
$2$-gerbe sur $B_G$ liée par $N$, munie d'une trivialisation au-dessus de
$X \simeq (B_G)_{/P}$. Ces $2$-gerbes forment a priori une $3$-catégorie de
Picard, mais ses $3$-flèches sont triviales, ce qui exprime
$H^0(B_G \bmod X, N) = 0$, c'est-à-dire l'injectivité de
$H^0(B_G, N) \to H^0(X, N)$ ; c'est donc une $2$-catégorie, et c'est celle des
Gr-champs épinglés.\footnote{La page écrit $H^0(B_G/X, N)$ pour
$H^0(B_G \bmod X, N)$. Le mot « un » de « cela un fait qui exprime » est lu
avec doute.}

Sur le topos ponctuel, la description redonne le théorème de Sính de la
page 3. Alors $B_G$ est le topos classifiant du groupe $G$ et
$R\Gamma(B_G, N)$ calcule la cohomologie de groupe. La suite exacte de la
cohomologie relative donne $H^{i}(B_G \bmod \mathrm{pt}, N) = H^{i}(G, N)$
pour $i \geq 2$. Les classes de Gr-catégories épinglées par $(G, N)$, qui sont
le $H^2$ du complexe, forment donc $H^3(G, N)$. Les auto-équivalences
épinglées à isomorphisme près, son $H^1$, forment $H^2(G, N)$. Les
automorphismes de l'identité, son $H^0$, forment le groupe
$H^1(B_G \bmod \mathrm{pt}, N) = Z^1(G, N)$ des
dérivations.\footnote{Cette vérification est la nôtre.}

Pour localiser sur $X$ et décrire le $2$-champ de Picard des Gr-champs
épinglés par $(G, N)$ sur des objets variables de $X$, on prend
\[ \tau_{\leq 2}\Bigl( R p_{G*}\, \mathrm{Coker}\bigl(N \to q_{G*}\,
   C(q_G^{*} N)\bigr) \Bigr) , \]
où $C$ désigne une résolution injective.\footnote{La page écrit « des champs
de Picard (sur des objets variables de $X$) épinglés par $G$, $N$ » ; il
s'agit des Gr-champs, $G$ n'étant pas supposé commutatif, et c'est de
Gr-champs que parle tout le reste de la page. « de Picard » est inséré après
« $2$-champ ». « rés.\ inj.\ » est écrit au-dessus du $C$. La page écrit
$R q_{G*}$ devant $C(q_G^*N)$ ; sur un complexe d'injectifs $R q_{G*}$ et
$q_{G*}$ coïncident.} C'est la version relative à $X$ du complexe global, la
construction a) de la section suivante appliquée à $q_G$ : le conoyau y
remplace la fibre décalée de $[1]$. Ces descriptions étant compatibles aux
variations de $G$, $N$, $X$, elles décrivent en principe la $2$-catégorie de
tous les Gr-champs.\footnote{Les $2$-gerbes ont été étudiées ensuite par
L.~Breen, \emph{On the classification of $2$-gerbes and $2$-stacks}
(Astérisque 225, 1994) ; la lettre ne pouvait pas le citer.}

\pagerange{12}{13}
\section*{La cohomologie relative (pages 12 et 13)}

Soient $q : X \to Y$ un morphisme de topos et $F$ un faisceau abélien, ou un
complexe, sur $Y$. Peut-on définir $R\Gamma(Y \bmod X, F)$
\emph{fonctoriellement} en $F$, de la catégorie dérivée de $\mathrm{Ab}(Y)$
dans celle de $\mathrm{Ab}$ ? L'interprétation par des $n$-champs de Picard
suggère que oui. Grothendieck ne voit de construction à la main que dans deux
cas extrêmes. Notons $C$ une résolution injective.\footnote{« résol.\ inj.\ »
est écrit sous le $C$ à la page 12, « rés.\ inj.\ » sous le premier $C$ à la
page 13 ; l'indice $Y$ de $R\Gamma_Y$ est empâté.}

\begin{enumerate}
\item[a)] $q$ est \emph{$(-1)$-acyclique} : $F \to q_* q^* F$ est injectif
pour tout $F$. C'est le cas de $Y_{/P} \to Y$ lorsque $P \to e_Y$ est un
épimorphisme, donc de $B_e \to B_G$ à la section précédente.\footnote{Le
renvoi « plus haut » est lu avec doute ; c'est lui qui date la page après les
pages 14 et 15.} On prend
$R\Gamma_Y\bigl(\mathrm{Coker}(F \to q_* C(q^*F))[-1]\bigr)$ : la flèche
étant injective en chaque degré, son conoyau est quasi-isomorphe au cône, et
le conoyau décalé de $-1$ calcule la fibre de
$R\Gamma(Y, F) \to R\Gamma(X, q^*F)$.
\item[b)] Pour tout $F$ injectif, $q^* F$ est injectif et $F \to q_* q^* F$
est un épimorphisme ; c'est le cas de l'inclusion d'un ouvert
$U \subset e_Y$. On prend
$R\Gamma_Y\bigl(\mathrm{Ker}(C(F) \to q_* q^* C(F))\bigr)$, le noyau d'une
flèche surjective en chaque degré calculant la fibre.\footnote{Après
$q_* q^*$ le trait d'ouverture est fait comme un $C$, et l'on pourrait lire
$q_* q^* C(C(F))$ ; on lit $q_* q^* C(F)$, seule lecture où la formule
calcule ce qu'elle doit.}
\end{enumerate}

Dans le cas général, la difficulté est que le cône d'un morphisme n'est pas
fonctoriel, dans la catégorie dérivée, par rapport à la flèche dont on prend
le cône. « Et pourtant, dans les cas particuliers actuels, il devrait y avoir
un choix fonctoriel. Est-ce évident ? »\footnote{Un mot est biffé dans la
parenthèse qui donne l'exemple $F \to q_*(q^*(F))$.} La réponse est oui, et
pour une raison générale : la flèche $F \to Rq_* q^* F$ n'est pas une flèche
quelconque de la catégorie dérivée, c'est l'unité d'une adjonction, et elle se
relève au niveau des complexes. Une catégorie de Grothendieck admet des
résolutions injectives fonctorielles ; avec elles $F \mapsto
\bigl(C(F) \to q_* C(q^* C(F))\bigr)$ est une transformation naturelle de
foncteurs de complexes, son cône est fonctoriel au niveau des complexes et
respecte les quasi-isomorphismes, et il descend donc à la catégorie
dérivée.\footnote{Cette réponse est la nôtre. Dans le langage d'aujourd'hui,
la catégorie dérivée est le quotient homotopique d'une $\infty$-catégorie
stable, où la fibre d'une transformation naturelle est fonctorielle ;
l'obstacle que la lettre nomme est propre aux catégories triangulées. Les
deux constructions a) et b) sont des cas où le cône se calcule par un
conoyau ou un noyau, et pour cette raison n'exigent rien de plus.}

\emph{Question pour Illusie.} Dans sa théorie des déformations de schémas en
groupes plats, Illusie rencontre des groupes $H^3(B_G \bmod X, \ \cdot\ )$ et
$\mathrm{Ext}^2(X, \ \cdot\ )$. Peut-on court-circuiter cette théorie par celle,
« supposée écrite », des Gr-champs --- ou par celle des champs de Picard ?
La phrase qui suit commence par « J'ai » et le dossier n'en contient pas la
suite.\footnote{Les seconds arguments sont laissés en blanc sur la page ; la
lettre dans le $\mathrm{Ext}^2$ est douteuse ($X$ ou $K$, avec un point
souscrit). La page écrit $H^3(B_G/X, \ )$ ; on lit la barre comme à la
page 15, pour la cohomologie relative. La théorie visée est celle de
L.~Illusie, \emph{Complexe cotangent et déformations} II, chapitre VII
(Lecture Notes in Math.\ 283, 1972) ; on n'identifie pas ici les groupes exacts
qu'elle emploie.}

\pagerange{16}{17}
\section*{Lettre à Deligne, Giraud et Verdier du 23.6.1974 (pages 16 et 17)}

Lettre tapée, datée de Villecun, signée à la main « Schurik ». Grothendieck
y présente le travail de Hoàng Xuân Sính, mené depuis son séjour au Vietnam
en 1967 sur un sujet qu'il lui avait proposé : une théorie de structure des
Gr-catégories et la construction d'une Gr-catégorie, ou plutôt d'une
catégorie de Picard, enveloppante pour une catégorie munie d'un produit
$\otimes$ associatif et commutatif --- c'est-à-dire, exactement, les deux
moitiés du tapuscrit des pages 1 à 6. Il a reçu une copie du manuscrit au
propre, obtenue « au prix de mille difficultés », et a adressé un rapport
officiel au ministre de l'enseignement supérieur et technique de la
République démocratique du Vietnam, dont il joint une copie.\footnote{Le
rapport annoncé « ci-joint » n'est pas dans le dossier.} Sính, avec son
accord, a demandé à soutenir en France ; l'accord des autorités
vietnamiennes n'est pas encore acquis. Si elle figure, comme il l'a appris,
dans la délégation vietnamienne au Congrès international de Vancouver
d'août 1974, elle pourrait soutenir à Orsay au retour, fin septembre ou début
octobre.

Il leur demande donc s'ils accepteraient en principe de faire partie du jury,
avec lui, et propose d'envoyer à qui le demande soit la thèse complète, soit
un résumé assez détaillé rédigé par Sính, soit un résumé moins détaillé « de
13 pages manuscrites » rédigé en vue des Comptes rendus mais trop long pour
eux ; à défaut, de chercher autour d'eux des mathématiciens « dans le bain »
du genre de choses que fait Sính.\footnote{« vous vous trouvez » et « dans le
bain » du genre de choses sont vérifiés sur le fac-similé. Rien dans le
dossier ne dit lequel de ces textes est le tapuscrit des pages 1 à 6 ; le
manuscrit anglais de treize pages du dossier 135 répond par sa longueur au
troisième, sans que le dossier présent permette de l'établir.} La page 17
confie les formalités à un frère de Sính établi en région parisienne, dont
l'adresse n'est pas reproduite,\footnote{La transcription ne reproduit pas
la ligne qui donne l'adresse postale et le téléphone : ce sont les données
d'une personne privée.} et un post-scriptum demande à Verdier de transmettre
la lettre et le rapport à Giraud, dont Grothendieck n'a pas l'adresse,
ajoutant que ce rapport pourra servir aussi de rapport de thèse auprès de
l'administration universitaire française.

\end{document}
